% SECCIÓN 4: EL CONSENTIMIENTO INFORMADO — CONFIGURACIÓN JURÍDICA Y EFECTOS DEL INCUMPLIMIENTO
%
% Objetivos de esta sección:
%   1. Caracterizar el consentimiento informado (CI) como derecho fundamental del paciente
%      y deber jurídico del médico (art. 15 Ley 26842; arts. 4, 8 Ley 29414).
%   2. Analizar los requisitos del CI válido: información suficiente, comprensible,
%      oportuna; voluntad libre e informada; forma (escrita vs. verbal).
%   3. Distinguir el CI como:
%      - Eximente de responsabilidad (cuando es válido y previo al acto médico).
%      - Fuente autónoma de responsabilidad (cuando hay déficit informativo que
%        causa daño, aunque el acto médico sea técnicamente correcto).
%   4. Analizar el nexo causal en el CI: ¿habría el paciente rechazado el tratamiento
%      de haber sido informado adecuadamente? (test de causalidad hipotética).
%   5. Problemas especiales: urgencia médica, incapacidad del paciente, CI por representación.
%
% Referencias clave: Ley 29414 (arts. 4, 8, 15), Fernández Cruz, Galán Cortés
% Extensión estimada: 1.500–2.000 palabras

\section{El consentimiento informado: configuración jurídica y efectos del incumplimiento}
\label{sec:consentimiento}

